\documentclass[../main.tex]{subfiles}

\begin{document}



% =======================================================================
\chapter{Experimental Results and Conclusions}
%\addcontentsline{toc}{chapter}{Experimental Results and Conclusions}
\label{chap:Empiric-analysis-models}
% =======================================================================

The previous chapter described the full experimental grid, from which the grand total of $\approx 13$k single models were trained (summarized in Section \ref{sec:experiment-protocol}). The present chapter reports and interprets the \emph{quantitative outcomes} of that grid. In line with the analytical priorities stated in the introduction of Chapter \ref{chap:5-Empiric-analysis-models}, the discussion opens with \red{Experiment B} (the principal, full-dataset campaign run over 20 random seeds), because it is the broader and more controlled study, and only afterwards turns to \red{Experiment A} (the proof of concept experiment), which is best read as one concrete instantiation of the more general phenomena that \red{Experiment~B} characterises in depth. The chapter then closes with the overall conclusions of the thesis.

Throughout, the \textbf{Matthews Correlation Coefficient (MCC)} is the primary metric, reported separately for the binary task (\texttt{label\_generic}: benign vs.\ malicious) and the multi-class task (\texttt{Label}: attack subtype). MCC is preferred over accuracy and macro-F1 because, as motivated in Section \ref{sec:hpo}, the CIC-IDS2018 label distribution is heavily imbalanced: MCC remains informative under imbalance and is bounded in $[-1, +1]$, where $0$ denotes a prediction uncorrelated with the ground truth and negative values denote systematic anti-correlation; yet, in contrast, unlike the F1-score, the MCC is much less susceptible to the presence of underrepresented classes within the target columns (which, as we know, are very common in the dataset used). For this reason, in this section, most of the analysis will be based on the MCC (although all F1-score thresholds have still been traced). Also, the aggregated F1 and accuracy tables produced by the evaluation notebooks are consistent with the MCC trends discussed here and are therefore not reproduced in full.

% =======================================================================
\newpage
\section{Reading the Results: Evaluation Protocol}
\label{sec:res-protocol}
% =======================================================================

Before presenting the numbers, it is useful to fix precisely how each reported quantity was computed, since every claim in this chapter is expressed in these terms.

\paragraph{Performance delta against the clean baseline.}
For a given model, task, corruption method, target column (or column group), and noise percentage $perc$, the metric of interest is the \emph{signed difference} between the mean MCC (indicated as $\overline{\mathrm{MCC}}$) of the noisy configuration  and the MCC of the corresponding clean baseline of Section \ref{sec:baseline-performance}:
\[
    \Delta\mathrm{MCC} \;=\; \overline{\mathrm{MCC}}_{\text{noisy}}(perc) \;-\; \mathrm{MCC}_{\text{baseline}}.
\]
A \emph{positive} $\Delta\mathrm{MCC}$ therefore means that injecting noise into the training data \emph{improved} held-out performance relative to training on clean data; a negative value denotes degradation. The clean validation and test partitions are identical across all configurations (Section~\ref{sec:noise-construction}), so these deltas are directly comparable across methods, features, and seeds.

\paragraph{Confidence intervals: the Student-$t$ interval.}
Each cell of the grid is repeated over the seed set of its phase (5 seeds for \red{Experiment~A}, 20 for \red{Experiment~B}), yielding, for every configuration, a small sample of $n$ MCC measurements. The reported figure is their mean $\mathrm{MCC}$ (indicated as $\overline{\mathrm{MCC}}$); but a mean over this number of runs could still be an uncertain estimate of the ``true'' expected performance, and that uncertainty must be quantified. Since the population variance is unknown and is estimated from the same handful of seeds, the appropriate tool is the \emph{Student-$t$ confidence interval}, reported throughout as $\overline{\mathrm{MCC}} \pm \delta$ (and drawn as a shaded band in every figure), with half-width
\[
    \delta \;=\; t_{1-\alpha/2,\;n-1}\;\frac{s}{\sqrt{n}},
\]
where:
\begin{itemize}
    \item $s$ is the sample standard deviation across seeds
    \item $s/\sqrt{n}$ is the standard error of the mean
    \item $t_{1-\alpha/2,\,n-1}$ is the critical value of Student's $t$ distribution with $n-1$ degrees of freedom at confidence level $1-\alpha = 0.95$
\end{itemize}

Student's $t$ (rather than the Gaussian $z$) is used precisely because $\sigma$ is itself estimated from few points: its heavier tails widen the interval to account honestly for that extra uncertainty. The interpretation is the customary one: were the experiment repeated with fresh seeds, approximately $95\%$ of the intervals so constructed would contain the true mean. The half-width therefore \emph{adapts to the seed count} through two compounding factors:
\begin{itemize}
    \item the $\sqrt{n}$ in the denominator
    \item the critical value $t_{1-\alpha/2,\,n-1}$, which shrinks from $2.776$ at $n=5$ (\red{Experiment~A}) to $2.093$ at $n=20$ (\red{Experiment~B}) and approaches the Gaussian limit of $1.960$ as $n\!\to\!\infty$, so that the 20-seed \red{Experiment~B} intervals are intrinsically tighter than the 5-seed \red{Experiment~A} ones
\end{itemize}

Therefore the deltas of this chapter are reported as: $$\Delta\mathrm{MCC} \pm \delta$$.

\paragraph{Statistical significance: the meaning of $p$.}
A narrow interval of $\delta$ tells us the mean is well estimated, but to decide whether an observed delta is \emph{real}, rather than an artifact of a particular random split or noise realization. For this reason the evaluation pipeline additionally runs a paired \emph{Wilcoxon signed-rank test} on the per-seed deltas of each noisy cell against the clean baseline, computed on the binary MCC (the most stable of the two tasks and the one most relevant to the anomaly-detection use case). The resulting \emph{$p$-value} is the probability of observing a sign-consistent effect at least this strong \emph{if noise truly had no effect} (the null hypothesis). A small $p$ therefore means that such a coherent shift across independent seeds would be exceedingly unlikely to arise by chance, and when: $$p < 0.05$$the null hypothesis is rejected and the effect flagged as significant. This is closely related to the confidence interval: an effect is significant, roughly, when its interval $\Delta\mathrm{MCC} \pm \delta$ lies entirely on one side of zero, so the intuitive reading of ``$p < 0.05$'' is ``the $\pm\delta$ margin is small enough, relative to the size of the effect, that we are confident of its sign''. The stronger value $p < 10^{-4}$ reported for several \red{Experiment~B} cells corresponds to all twenty seeds agreeing on the sign of the effect, whose probability under the null is $\approx 2\times 10^{-6}$: about as close to certainty as twenty seeds permit. One caveat follows directly from this construction, however: the $p$-value depends not only on the effect size and its $\pm\delta$ margin but also on the \emph{number of seeds}. With only five seeds the Wilcoxon test cannot return a value below $\approx 0.0625$ even when every seed agrees, which is why \red{Experiment~A} can display large, tightly-bounded effects that the formal test is nevertheless unable to certify (discussed in Section~\ref{sec:res-A}).

\paragraph{The sample-size asymmetry between models.}
As established in Section~\ref{sec:sample-dimensionality}, TabNet evaluates on roughly ten times as many samples as CNN-LSTM (one flow per sample versus one sliding window of $50$ flows per sample). Consequently, CNN-LSTM metrics are intrinsically noisier and exhibit wider seed-to-seed variance; the 20-seed confidence intervals of \red{Experiment~B} are precisely what make it possible to separate genuine effects from this elevated variance. This asymmetry should be kept in mind whenever the magnitude of a CNN-LSTM delta is compared against a TabNet one.

% =======================================================================
\newpage
\section{Experiment B: Results on full theorethical analysis}
\label{sec:res-B}
% =======================================================================

\red{Experiment~B} is the analytical core of the thesis. It operates on the complete (non feature-pruned) dataset, restricted at training time to the nine-feature working set of Table~\ref{tab:final-features} plus the \texttt{Timestamp} column, and is repeated over the 20 seeds $S_B$ with the six-level grid $p_B \in \{5\%, 10\%, 20\%, 35\%, 50\%, 75\%\}$. Its purpose is not to maximise detection rates (AKA improving the NIDS system) but to map, as completely as possible, the space of effects that controlled corruption produces on a tabular classifier (TabNet) and on a time-series classifier (CNN-LSTM). 
The reference points for this Experiment are the clean baselines of Tables~\ref{tab:tabnet-baseline} and~\ref{tab:cnnlstm-baseline}, namely:
\begin{itemize}
    \item For TabNet:
    \begin{itemize}
        \item  $\mathrm{MCC}_{\text{bin}} = 0.883$
        \item  $\mathrm{MCC}_{\text{mul}} = 0.856$
    \end{itemize} 
    \item For CNN-LSTM:
    \begin{itemize}
        \item  $\mathrm{MCC}_{\text{bin}} = 0.737$
        \item  $\mathrm{MCC}_{\text{mul}} = 0.635$
    \end{itemize} 
\end{itemize}

The campaign is divided into:
\begin{itemize}
    \item \emph{single-feature} scenario (one column corrupted at a time, Section~\ref{sec:res-B-single})
    \item \emph{multiple-feature} scenario (several columns corrupted jointly, Section~\ref{sec:res-B-multi})
\end{itemize}
This is done in order to mirror the grid definitions of Sections \ref{sec:exp-B-single} and \ref{sec:exp-B-multi} from previous Chapter.

% ------------------------------------------------------------
\begin{landscape}
\subsection{Single-Feature Corruption}
\label{sec:res-B-single}
% ------------------------------------------------------------
Full display of all conducted experiments using \texttt{Pucktrick} Library is displayed in images \ref{fig:TabNet_full_grid_single_feature} and \ref{fig:CNN-LSTM_full_grid_single_feature} below: 
\begin{figure}[H]
    \centering
    \includegraphics[width=1\linewidth]{images/results/TabNet_full_grid_single_feature.png}
    \caption{Full result grid spectrum over MCC for TabNet in both prediction heads (bynary task sx, multitask task dx)}
    \label{fig:TabNet_full_grid_single_feature}
\end{figure}
\end{landscape}

\begin{landscape}
\begin{figure}[H]
    \centering
    \includegraphics[width=1\linewidth]{images/results/CNN-LSTM_full_grid_single_feature.png}
    \caption{Full result grid spectrum over MCC for CNN-LSTM in both prediction heads (bynary task sx, multitask task dx)}
    \label{fig:CNN-LSTM_full_grid_single_feature}
\end{figure}
\end{landscape}


% ------------------------------------------------------------
\newpage
\subsubsection{Corruption of input features}
\label{sec:res-B-single-feat}
% ------------------------------------------------------------

The single-feature feature-corruption sweep applies the three feature-level \texttt{Pucktrick} methods (\texttt{missing}, \texttt{noise}, \texttt{outliers} and the row-level \texttt{duplicated} method to each of the ten working columns. Three qualitatively distinct regimes emerge, and they recur throughout the entire campaign.

\paragraph{Regime 1: Both Models are mostly resilient against single feature noise injection.}
As we can see from Images \ref{fig:TabNet_full_grid_single_feature} and \ref{fig:CNN-LSTM_full_grid_single_feature}, for both models, corrupting any single input feature leaves the binary MCC, while decreased, essentially unchanged across: \emph{most} of the ten-feature $\times$ five \texttt{Pucktrick} methods $\times$ six-percentage grid. Especially for TabNet, which is more robust to feature corruption. This robustness is the empirical signature of TabNet's sequential attention mechanism: when an informative feature is corrupted, the attentive transformer simply reallocates its sparse attention mass towards the remaining clean features, and because the working set deliberately contains redundant discriminative signal, binary detection barely suffers. The multi-class task shows a small but systematic offset (a roughly uniform $-0.02$ to $-0.03$ penalty under almost any corruption, including corruption of zero-variance control features), confirming that the resolution of rare attack subtypes is the first capability to erode when the input distribution is perturbed, even when the benign/malicious boundary is untouched.

If we take a look at the binary-label portion of those graphs: it is still observable in the majority of cases the tendency to decrease to a lower MCC than the baseline, the binary $\Delta\mathrm{MCC}$ remains within $\pm 0.1$-$\pm 0.4$ in the vast majority of cells. Yet, there is an exception to this. The only feature whose corruption produces a sizable, monotone degradation for both models is \texttt{Dst Port}; which represents:
\begin{itemize}
    \item For \textbf{CNN-LSTM}: the \emph{second} most discriminative column (as referenced in Table \ref{tab: final-ranking-features-TabNet-baseline})
    \item For \textbf{TabNet}: the \emph{first} most discriminative column (as referenced in Table \ref{tab:final-ranking-features-CNN-LSTM-baseline})
\end{itemize}

\begin{figure}[H]
    \centering
    \includegraphics[width=0.6\linewidth]{images/results/DTS-Port-comparaison_singleFeature.png}
    \caption{Most evident case of performance degradation on Experiment-B (single feature noise injection): happening on feature `DTS Port`}
    \label{fig:CNN-single-feature-dts-port}
\end{figure}

At first glance, while it is completely reasonable for TabNet that the most discriminative feature when injected whit noise might report the most damage, it might seem strange that only the \emph{second} most important feature for CNN-LSTM creates this same distortion to MCC value, while the \emph{first} most important feature doesn't apply the same level of reduction to the MCC ratio. Yet, by taking a deeper look into the data, this behavior is completely explicable.
If we take a look at CNN-LSTM most important feature it results being: $$\texttt{Fwd Seg Size Min}$$
This particular feature, is sole one that we identified as "Low-Fragility" inside the whole dataset inside Dataset Analysis Chapter \ref{sec:fragility} (Specifically at Table \ref{tab:fragility-examples}), justifying why at low levels of noise injections (below 35\%) no evident results of declining MCC were shown, and only reaching a 4\% decline of MCC at 70\% noise. 

This statement could further be backed up by taking a look at following graph: 

\begin{figure}[H]
    \centering
    \includegraphics[width=0.8\linewidth]{images/results/top_feature_ranking_evolution.png}
    \caption{Average (over 20 random seeds) position ranking taken by top-2 features for both TabNet and CNN-LSTM across increasing noise injection.}
    \label{fig:top_feature_ranking_evolution}
\end{figure}

\newpage
Graph \ref{fig:top_feature_ranking_evolution} shows that:  
\begin{itemize}
    \item For \textbf{CNN-LSTM}: 
    \begin{itemize}
        \item For Feature \texttt{Fwd Seg Size Min}: This result further consolidates the notion that feature \texttt{Fwd Seg Size Min} is highly resistant to noise injection methods, as for all noise injection experiments up to 35\% noise the feature always remains the in the average number one spot. After the 35\% experiment (at 50\%) the feature average rank drop significantly (for all examined noise injection methods). By taking a look at Image \ref{fig:CNN-LSTM_full_grid_single_feature} (at this particular feature graph), the performance actually start drastically declining at the same point. 
        \item For feature \texttt{DTS Port}: we can verify that this feature immediately drops position to 3rd and subsequently to 5th after first noise injection experiment, further backing up results in Image \ref{fig:CNN-single-feature-dts-port}, where it is shown (for all noise injection methods) immediatly causing MCC to drop. 
    \end{itemize}
    \item For \textbf{Tabnet}: We can also observe that feature \texttt{DTS Port} immediately drops after just 5\% noise injection for all methods, yet since this column.
\end{itemize}

To summarize the points discussed in this subsection, it appears that both models are particularly resistant to noise injection (especially TabNet, thanks to its sequential attention mechanism): a resilience reflected in the narrow Student-$t$ confidence intervals around these near-baseline deltas. At the same time, when noise is introduced into a feature that is particularly important for a model, and the model no longer has redundant features to rely upon to reach the correct solution, performance inevitably declines --- a degradation that the paired Wilcoxon test confirms as highly significant ($p < 10^{-4}$, i.e.\ all twenty seeds agree on the loss): this point will be reconfirmed in the next section, which analyzes the experiments conducted on groups of features (\ref{sec:res-B-multi}).



\paragraph{Regime 2: CNN-LSTM \emph{benefits} from corrupting the \texttt{Timestamp} feature.}
The most striking single-feature result is that injecting noise into the \texttt{Timestamp} column \emph{improves} the CNN-LSTM on both tasks, and the improvement is large, monotone in the noise level, and highly significant. Table \ref{tab:res-B-timestamp} reports the effect for the \texttt{outliers} method, which is the strongest case; \texttt{missing} and \texttt{noise} on \texttt{Timestamp} produce smaller but equally consistent gains (binary $\Delta\mathrm{MCC} \approx +0.04$ to $+0.05$ across all percentages for \texttt{missing}). All \texttt{outliers}/\texttt{Timestamp} cells from $10\%$ upward are significant against the clean baseline having $p$-value of: $$p < 10^{-4}$$which provides a \emph{the highest} degree of certainty registered across all experiments.

\begin{table}[H]
\centering
\caption{\red{Experiment~B}, single-feature: CNN-LSTM MCC improvement under \texttt{outliers} corruption of the \texttt{Timestamp} feature. Each entry is $\Delta\mathrm{MCC} \pm \delta$ against the clean CNN-LSTM baseline ($0.737$ binary, $0.635$ multi-class), where $\delta$ is the Student-$t$ $95\%$ confidence half-width over the 20 seeds. The margins ($\le 0.024$) are far smaller than the gains, which is precisely why all cells from $10\%$ upward are significant at $p < 10^{-4}$.}
\label{tab:res-B-timestamp}
\small
\resizebox{\linewidth}{!}{%
\begin{tabular}{lrrrrrr}
\toprule
\textbf{Task} & \textbf{5\%} & \textbf{10\%} & \textbf{20\%} & \textbf{35\%} & \textbf{50\%} & \textbf{75\%} \\
\midrule
Binary       & $+0.037{\pm}0.010$ & $+0.068{\pm}0.013$ & $+0.081{\pm}0.009$ & $+0.110{\pm}0.012$ & $+0.147{\pm}0.009$ & $+0.139{\pm}0.004$ \\
Multi-class  & $+0.082{\pm}0.022$ & $+0.123{\pm}0.024$ & $+0.164{\pm}0.012$ & $+0.206{\pm}0.014$ & $+0.233{\pm}0.010$ & $+0.236{\pm}0.006$ \\
\bottomrule
\end{tabular}%
}
\end{table}

\begin{figure}[H]
    \centering
    \includegraphics[width=1\linewidth]{images/results/Timestamp_comparaison_singleFeature.png}
    \caption{\red{Experiment~B}, single-feature: CNN-LSTM held-out MCC as a function of the \texttt{Timestamp} corruption level, one line per PuckTrick method, with Student-$t$ confidence bands over the 20 seeds. The dashed horizontal line is the clean baseline.}
    \label{fig:res-B-timestamp}
\end{figure}

This phenomenon gives us a very strong result: corrupting the temporal index (particularly with \texttt{outliers}) acts as a regularizer and raises performance well above the clean reference for both tasks (single and multi-label).

In fact, although the CNN-LSTM model operates on chronologically ordered sliding windows (Section~\ref{Chapter: model cnn-lstm}), the raw \texttt{Timestamp} value is \emph{never} presented to the network as an input channel. In the preprocessing pipeline the column is retained only transiently: it is used as the sort key of the temporal \texttt{orderBy} operation that arranges the flows in chronological order, and it is subsequently \emph{dropped} before the feature tensor is materialized. Consequently, the model cannot predict a label from \emph{when} a flow occurred, since the temporal coordinate is not part of its observable input.

This does not, however, imply that the temporal dimension is inert. Because the windows are built over a chronologically sorted sequence and attack campaigns are clustered in time, the labels of flows that fall within the same window are strongly auto-correlated. The CNN-LSTM can therefore exploit a \emph{temporal shortcut}: rather than discriminating each flow solely from its own statistical signature, it can lean on
the local label homogeneity of its neighbors within the window (i.e.\ "the surrounding flows are all of attack type $X$, hence this one is likely $X$ as well"). Crucially, this shortcut is mediated by the \emph{ordering-and-windowing} mechanism.

This observation explains a mechanism through which noise injection on the \texttt{Timestamp} column can influence the CNN-LSTM in a manner that is \emph{indirect}, yet substantial. In the experimental pipeline, corruption is applied to the training partition \emph{before} the chronological sort is performed. When the corrupted column is \texttt{Timestamp} itself (for the value-perturbing methods \texttt{missing}, \texttt{noise} and \texttt{outliers}), the subsequent \texttt{orderBy} no longer reflects the true arrival order: the rows are effectively \emph{permuted}, and the sliding windows are assembled over a partially shuffled sequence. The corruption
therefore does not pollute the feature values consumed by the model; instead, it \emph{reshuffles the samples prior to window construction}, weakening the within-window label auto-correlation on which the temporal shortcut relies.

Therefore: the effect is best understood as a form of \emph{regularisation}: by disrupting the spurious temporal clustering in the training windows, the corruption forces the network to rely more heavily on the genuine flow-level statistical signatures rather than on the contextual label homogeneity of adjacent flows. This pressure is further accentuated by the asymmetry of the experimental protocol: noise is injected \emph{only} into the training partition, whereas the validation and test partitions retain their true chronological ordering. The model is thus trained on windows whose temporal shortcut has been degraded, but is evaluated on cleanly ordered windows. Under this regime, an
\emph{improvement} in CNN-LSTM performance when the \texttt{Timestamp} column is corrupted is consistent with the model being discouraged from over-fitting to the temporal clustering of attack campaigns and instead learning more transferable, feature-driven representations.

The effect is absent for TabNet precisely because TabNet, by contrast, operates on a \emph{two-dimensional} design matrix $(n_{\text{samples}}, n_{\text{features}})$ in which each flow is an independent row and no temporal window is ever formed. Since the model treats samples as exchangeable and the training batches are randomly shuffled regardless, the row ordering carries no information that TabNet can exploit. Corrupting the \texttt{Timestamp} column (which is dropped before the feature matrix is built) therefore leaves the set of feature vectors seen by TabNet effectively unchanged (it merely reorders identical rows). With no temporal shortcut to exploit, there is correspondingly none to regularize, and the manipulation has no material effect on TabNet's performance.




\paragraph{Regime 3: Row-level method \texttt{duplicated} has no effect on performance.}

The method \texttt{duplicated} was used on `\texttt{Timestamp}` only as a nominal column label, while the operation took place operating row-wise. This method's are therefore visible in Image \ref{fig:res-B-timestamp}. As we can see the effects are near-null on the MCC of both models (binary $\Delta\mathrm{MCC}$ within inter-seed noise across the whole grid). This may appear as counter-intuitive, since the duplication of training rows is sometimes regarded in the literature as a form of implicit regularization. The result is, however, exactly what the mechanics of the operation predict. Because \texttt{duplicated} inserts \emph{byte-identical} copies into the training partition only, it adds no entropy: an exact duplicate yields an identical forward pass and an identical gradient, so its effect is a mere \emph{re-weighting} of the affected rows rather than the injection of new information. Genuine regularization instead acts in one of three ways:
\begin{itemize}
    \item by penalizing model complexity (e.g.\ $L_2$),
    \item by introducing stochastic perturbation (e.g.\ dropout, label smoothing, \textsc{mixup}),
    \item or by expanding the data manifold with \emph{novel} points (data augmentation).
\end{itemize}
Exact duplication performs none of these. Formally, duplicating a uniformly sampled fraction $p$ of rows rescales the empirical distribution by an almost-constant factor that is absorbed by the learning rate and the per-epoch sample count, leaving the
empirical-risk minimizer, and hence the clean-set MCC, essentially invariant.
Three properties of the present setting reinforce this inertia: 
\begin{itemize}
    \item stochastic gradient descent already revisits every sample many times across epochs, so one additional copy of $p\%$ of the rows is marginal;
    \item the training partition is large ($\approx 1.07$ million flows), so duplicate draws from the same distribution do not perturb the gradient statistics;
    \item critically, the duplication is uniform rather than class-targeted, so it does not correct the label imbalance: the one mechanism through which row duplication is genuinely beneficial (minority-class oversampling)
\end{itemize}

The recurrent intuition that duplication regularities is, moreover, traceable to a conflation with \emph{bootstrap aggregating}, whose regularizing power stems from averaging an ensemble of decorrelated models, not from the duplicates contained within any single bootstrap sample. Consistently with this analysis, the empirical deltas hover within the seed-to-seed noise, with only a mild \emph{negative} drift on the CNN-LSTM multi-class head at high duplication levels ($-0.134$ at $20\%$, $-0.118$ at $75\%$): since duplicated rows share the \texttt{Timestamp} of their source and therefore fall adjacent to it once the sequence is chronologically ordered, a sliding window may contain repeated identical flows, slightly distorting the local temporal statistics on which the CNN-LSTM relies: an effect that, if anything, mildly harms rather than helps, and which leaves the order-invariant TabNet untouched. 

There is also an in-domain explanation to this effect: as we said before in Chapter \ref{sec:temporal-split}. Inside network-flows it is common that duplicate packets flow though the network, adding some more duplicates to the dateset practically means keeping intact the same dataset structure that was already present as some dataset rows were already inevitably duplicated. Therefore makes perfect sense that no effects are recorded using this particular method, in this particular case study.



% ------------------------------------------------------------
\newpage
\subsubsection{Corruption of the target labels}
\label{sec:res-B-single-label}
% ------------------------------------------------------------

A complementary script applies the \texttt{labels} \texttt{Pucktrick} method directly to both \emph{target} columns. This way we aim to isolate the effect of corrupting the supervision signal rather than an input feature. The two targets are corrupted \emph{independently}, one at a time. Crucially, because both architectures are \emph{multi-task} (Section~\ref{Chap: model selection}), every corrupted target is evaluated on \emph{both} prediction heads at once. This is done because examining how the effects of corruption on one label impact the second label is considered a particularly interesting case study. Each of the two corrupted columns therefore yields two MCC curves per model:
\begin{itemize}
    \item the \emph{direct} effect, on the head that predicts the very column being corrupted
    \item  the \emph{cross} effect, on the \emph{other} head
\end{itemize}
for a total of $2\ \text{targets}\times 2\ \text{heads}\times 2\ \text{models} = 8$ series, all displayed in Figure~\ref{fig:res-B-labelcliff} and tabulated in Table~\ref{tab:res-B-labelcliff}. This cross-head view is one of the focal analyses of the thesis: it answers not merely ``how badly does corrupting a label damage its own task?'' but the more revealing question ``when one label is corrupted, what happens to the \emph{other} one?'': a question that only a multi-task probe can pose, and whose answer turns out to depend sharply on both \emph{which} target is corrupted and on the model architecture.

\begin{figure}[H]
    \centering
    \makebox[\textwidth][c]{\includegraphics[width=1.2\linewidth]{images/results/label_method_noise_singleFeature.png}}
    \caption{\red{Experiment~B}, single-feature, \texttt{labels} method on the two target columns. Each panel fixes a model and a monitored MCC head (binary / multi-class); within each panel the two curves correspond to corrupting \emph{ONE} of the two label columns, so that the \emph{direct} effect (corrupted target $=$ monitored head) and the \emph{cross} effect to the other head can be read side by side. Bands are Student-$t$ $95\%$ intervals over the 20 seeds; the dashed line is the clean baseline.}
    \label{fig:res-B-labelcliff}
\end{figure}


% -------------------------------
\begin{landscape}
    
\begin{table}[H]
\centering
\caption{\red{Experiment~B}, single-feature, \texttt{labels} method: $\Delta\mathrm{MCC} \pm \delta$ (Student-$t$ $95\%$ half-width over 20 seeds) for the \emph{eight} target-corruption series. The table is split by the \emph{corrupted} target column; within each block every model is evaluated on both prediction heads. A "$\dagger$" marks the \emph{direct} head (the one being actively corrupted by \texttt{Pucktrick}); the unmarked rows are the \emph{cross} head, whose target was \emph{not} corrupted.}
\label{tab:res-B-labelcliff}
\small
\resizebox{\linewidth}{!}{%
\begin{tabular}{llrrrrrr}
\toprule
\textbf{Model} & \textbf{Monitored head} & \textbf{5\%} & \textbf{10\%} & \textbf{20\%} & \textbf{35\%} & \textbf{50\%} & \textbf{75\%} \\
\midrule
\multicolumn{8}{l}{\emph{Corrupting the binary target} \texttt{label\_generic}} \\
\midrule
TabNet   & Binary$^\dagger$    & $-0.007{\pm}0.004$ & $-0.005{\pm}0.003$ & $-0.009{\pm}0.003$ & $-0.016{\pm}0.004$ & $-0.956{\pm}0.041$ & $-1.642{\pm}0.035$ \\
TabNet   & Multi-class         & $-0.030{\pm}0.008$ & $-0.024{\pm}0.005$ & $-0.028{\pm}0.006$ & $-0.033{\pm}0.004$ & $-0.056{\pm}0.022$ & $-0.050{\pm}0.015$ \\
\cdashline{1-8}
CNN-LSTM & Binary$^\dagger$    & $+0.008{\pm}0.009$ & $+0.010{\pm}0.007$ & $-0.003{\pm}0.011$ & $-0.010{\pm}0.011$ & $-0.668{\pm}0.114$ & $-0.864{\pm}0.094$ \\
CNN-LSTM & Multi-class         & $+0.018{\pm}0.025$ & $+0.019{\pm}0.028$ & $-0.010{\pm}0.029$ & $-0.017{\pm}0.021$ & $-0.008{\pm}0.022$ & $-0.536{\pm}0.065$ \\
\midrule
\multicolumn{8}{l}{\emph{Corrupting the multi-class target} \texttt{Label}} \\
\midrule
TabNet   & Binary              & $-0.003{\pm}0.002$ & $-0.003{\pm}0.005$ & $-0.005{\pm}0.005$ & $-0.002{\pm}0.003$ & $-0.004{\pm}0.003$ & $-0.002{\pm}0.002$ \\
TabNet   & Multi-class$^\dagger$ & $-0.028{\pm}0.004$ & $-0.030{\pm}0.006$ & $-0.033{\pm}0.006$ & $-0.029{\pm}0.004$ & $-0.034{\pm}0.003$ & $-0.042{\pm}0.007$ \\
\cdashline{1-8}
CNN-LSTM & Binary              & $+0.004{\pm}0.008$ & $+0.011{\pm}0.007$ & $+0.012{\pm}0.007$ & $+0.001{\pm}0.009$ & $-0.029{\pm}0.044$ & $+0.000{\pm}0.011$ \\
CNN-LSTM & Multi-class$^\dagger$ & $+0.004{\pm}0.031$ & $+0.008{\pm}0.030$ & $+0.016{\pm}0.026$ & $-0.030{\pm}0.030$ & $-0.123{\pm}0.053$ & $-0.049{\pm}0.031$ \\
\bottomrule
\end{tabular}%
}
\end{table}

Both Table \ref{tab:res-B-labelcliff} and Image \ref{fig:res-B-labelcliff} reveal two pretty interesting truths for  models: 
\begin{enumerate}
    \item \textbf{corrupting the binary target (\texttt{label\_generic}) drags the multi-class-target head down with it, at very high corruption levels} (Although loss on multiclass-label is much more significant in CNN-LSTM rather then TabNet where it is barely noticeable). 
    \item \textbf{corrupting the multi-class-target (\texttt{Label}) degrades only its own head and leaves the binary head at baseline}.
\end{enumerate}
Both of those cases are further explored in paragraphs below.
\end{landscape}
% -------------------------------

\paragraph{Direct consequences on label injected with noise:}
Reading the directly-corrupted heads (the $\dagger$ rows of Table~\ref{tab:res-B-labelcliff}), the two targets behave in qualitatively opposite ways. Corrupting the \emph{binary} target (\texttt{label\_generic}) leaves its own (binary) head essentially untouched up to $35\%$ and then collapses it off a cliff at $50\%$ ($-0.956$ for TabNet, $-0.668$ for CNN-LSTM) and into outright inversion at $75\%$ (, i.e.\ a TabNet binary MCC of $\approx -0.76$). This is exactly what theory predicts for a two-state symmetric label-flip: with only two classes the \texttt{labels} method necessarily \emph{flips} every selected entry, so below $50\%$ the majority of the supervision is still correct, at $50\%$ the signal is information-free (the flipped and correct fractions balance), and beyond it the labels become a minority-inverted signal that the model dutifully learns in reverse, yielding a \emph{negative} clean-set MCC. The $50\%$ flip point is thus a genuine phase transition, not a gradual decline. Therefore the widening confidence margins at the cliff (CNN-LSTM binary $\pm0.114$ at $50\%$ versus $\pm0.011$ in the flat regime) quantify the seed-to-seed instability of a model balanced on the knife-edge of inversion. 
Corrupting the \emph{multi-class} target (\texttt{Label}), by contrast, never triggers a significant cliff on its own (multi-class) head: degradation is gradual and monotone (TabNet a steady $-0.03$ to $-0.04$; CNN-LSTM down to $-0.123$ at $50\%$). With fifteen classes a $p\%$ corruption \emph{reassigns} each affected label to one of many alternatives rather than guaranteeing an inversion, so no coherent inverse mapping can ever form. 


\paragraph{Indirect consequences on label injected with noise: Leakage of prediction error:}
Although both models seem to not register any significant performance loss before 35\% noise injection, but at 50\% they both collapse tremendously (TabNet: $\Delta MCC = -0.956{\pm}0.041$, while CNN-LSTM: $\Delta MCC = -0.668{\pm}0.114$) on binary label being corrupted; while registering no particular sign of deterioration on multiclass-label. Yet, after that point, reaching 75\% level of corruption also the multiclass-label prediction start tremendously suffer registering levels of:
\begin{itemize}
    \item For \textbf{TabNet}:
    \begin{itemize}
        \item $\Delta MCC = -1.642$ : for binary target where noise is being injected (flipping binary label).
        \item $\Delta MCC = -0.050$ : for multiclass target (clean).
    \end{itemize}
    This signifies that TabNet is particularly resilient even to this more aggressive noise injection strategies, as \textbf{the clean target} (multiclass) only loses about very few points of MCC, while target that has been corrupted (binary). This makes sense as TabNet feeds its two heads from a sparsity-regularised, attention-selected representation that stays partially disentangled per task, so a corrupted supervisory signal remains largely confined to its own head.
    
    \item For \textbf{CNN-LSTM}:
    \begin{itemize}
        \item $\Delta MCC = -0.864$ : for binary target where noise is being injected (flipping binary label).
        \item $\Delta MCC = -0.5360$ : for multiclass target (clean).
    \end{itemize}
    A much more interesting result, the catastrophic collapse of \texttt{label\_generic} \emph{propagates} to the multi-class head of the CNN-LSTM. This makes complete sense when considering CNN-LSTM architecture: the CNN-LSTM's two heads branch off a \emph{single} shared context vector, so the catastrophic gradient of the inverted binary loss contaminates the multi-class head as well. 
\end{itemize}
Therefore we can conclude that \textbf{the contamination is architecture-dependent}: the very same binary collapse that drags the CNN-LSTM's multi-class head down by $-0.536$ barely moves TabNet's ($-0.050$ at $75\%$).  

Quite the contrary, when considering effects of corrupting multi-class-target on binary class prediction. In those cases no significant consequences on MCC are monitored. Even if considering that CNN-LSTM shares prediction backbone for both prediction heads, we easily proved that damage flows from the binary target to the multi-class one, but not the reverse precisely because the binary collapse is a \emph{global inversion} of the representation that the two heads share, whereas multi-class corruption is a diffuse re-shuffle carrying no coherent signal to leak. 



% ------------------------------------------------------------
\begin{landscape}
\subsection{Multiple-Feature Corruption}
\label{sec:res-B-multi}
% ------------------------------------------------------------

The multiple-feature scenario corrupts several columns simultaneously within a single \texttt{PuckTrick} call, probing the compound effects that the single-feature sweep cannot reach. Six input-feature groups (spanning combinations of importance and fragility, Section~\ref{sec:exp-B-multi}) are corrupted with \texttt{missing}, \texttt{noise}, and \texttt{outliers}, while a seventh group bundles the two targets for joint \texttt{labels} corruption. Full experimental grid results are displayed below:
\begin{figure}[H]
    \centering
    \makebox[\textwidth][c]{\includegraphics[width=1.3\linewidth]{images/results/Expb_multifeature_CNN_full.png}}
    \caption{Full result grid spectrum over MCC for CNN-LSTM in both prediction heads (bynary task sx, multitask task dx)}
    \label{fig:xpb_multifeature_CNN_full}
\end{figure}
\end{landscape}

\begin{landscape}
\begin{figure}[H]
    \centering
    \makebox[\textwidth][c]{\includegraphics[width=1.3\linewidth]{images/results/Expb_multifeature_TabNet_full.png}}
    \caption{Full result grid spectrum over MCC for TabNet in both prediction heads (bynary task sx, multitask task dx)}
    \label{fig:Expb_multifeature_TabNet_full}
\end{figure}
\end{landscape}

% ------------------------------------------------------------
\subsubsection{Corruption of feature groups}
\label{sec:res-B-multi-feat}
% ------------------------------------------------------------

Compound corruption is, broadly, the additive-to-superadditive amplification of the single-feature degradations. Table~\ref{tab:res-B-multi-outliers} reports the degradation of \emph{both} models under \texttt{outliers} at the highest noise level, ordered to expose the dependence on group composition and, side by side, the order-of-magnitude gap between the two architectures.
Notably, Table below only reports results just from \texttt{Pucktrick}'s \texttt{outliers} method as all other two methods applied (\texttt{missing} and \texttt{duplicated}) report similar effects but in smaller scale, making \texttt{outliers} the most interesting case study.

\begin{table}[H]
\centering
\caption{\red{Experiment~B}, multiple-feature: $\Delta\mathrm{MCC} \pm \delta$ (Student-$t$ $95\%$ half-width over 20 seeds) at $75\%$ \texttt{outliers} corruption, by feature group, for \emph{both} models.}
\label{tab:res-B-multi-outliers}
\small
\resizebox{\linewidth}{!}{%
\begin{tabular}{lrrrr}
\toprule
& \multicolumn{2}{c}{\textbf{CNN-LSTM}} & \multicolumn{2}{c}{\textbf{TabNet}} \\
\cmidrule(lr){2-3}\cmidrule(lr){4-5}
\textbf{Corrupted feature group} & \textbf{Binary} & \textbf{Multi-class} & \textbf{Binary} & \textbf{Multi-class} \\
\midrule
Top-3 importance (\texttt{Fwd Act Data Pkts}, \texttt{Dst Port}, \texttt{Init Fwd Win Byts}) & $-0.207{\pm}0.030$ & $-0.367{\pm}0.062$ & $-0.039{\pm}0.018$ & $-0.054{\pm}0.015$ \\
Top + middle tiers (6 features)                                                              & $-0.209{\pm}0.011$ & $-0.526{\pm}0.080$ & $-0.026{\pm}0.009$ & $-0.051{\pm}0.012$ \\
Middle tier (\texttt{Protocol}, \texttt{PSH Flag Cnt}, \texttt{Fwd Seg Size Min})            & $-0.059{\pm}0.022$ & $-0.153{\pm}0.051$ & $-0.013{\pm}0.007$ & $-0.029{\pm}0.005$ \\
Highest fragility (\texttt{Fwd Act Data Pkts}, \texttt{Idle Mean}, \texttt{Init Fwd Win Byts}) & $-0.038{\pm}0.022$ & $-0.126{\pm}0.046$ & $-0.009{\pm}0.002$ & $-0.026{\pm}0.006$ \\
One-per-tier (\texttt{Fwd Act Data Pkts}, \texttt{PSH Flag Cnt}, \texttt{Idle Mean})         & $-0.014{\pm}0.009$ & $-0.038{\pm}0.033$ & $-0.005{\pm}0.003$ & $-0.026{\pm}0.006$ \\
Zero-variance control (\texttt{Bwd Byts/b Avg}, \texttt{Fwd Blk Rate Avg}, \texttt{Idle Mean}) & $-0.010{\pm}0.010$ & $-0.035{\pm}0.026$ & $-0.004{\pm}0.006$ & $-0.039{\pm}0.009$ \\
\bottomrule
\end{tabular}%
}
\end{table}

Notably, as visible from Table \ref{tab:res-B-multi-outliers},  no column-group features the \texttt{Timestamp} column. This is done because this experiment's final objective is expand on previous experiment about single-feature-noise-injection (Section \ref{sec:res-B-single-feat}) most destructive performance and doubling down on them, figuring out if any of this corruptions of multiple columns could reach  same levels of  prediction-performance-decay observed on corruption of binary label (seen in section \ref{sec:res-B-single-label}).
Therefore, the \textbf{beneficial regime essentially disappears in the compound setting}: the \texttt{Timestamp} feature is not a member of any group, so the regularizing effect of Section~\ref{sec:res-B-single-feat} cannot manifest, and only a handful of marginal CNN-LSTM gains survive (e.g.\ \texttt{noise} on the mixed-fragility group, $+0.010$ at $35\%$).

\paragraph{Corruption on Binary-Task classification head:}
It is also visible from Table \ref{tab:res-B-multi-outliers} that degradation scales jointly with the number of corrupted features and their importance. Yet the two architectures respond on entirely different scales: on the worst group (six important features) the CNN-LSTM binary-class head loses $\Delta MCC =-0.209$, whereas TabNet loses only $\Delta MCC = -0.026$, roughly an order of magnitude less.

This is a trend of all corruption experiments by feature-groups methods applied, as we can see from Image below:

\begin{landscape}
    
\begin{figure}[H]
    \centering
    \includegraphics[width=0.8\linewidth]{images/results/expB-multi-feature-full-point.png}
    \caption{Grouping of multi-feature experiments by \texttt{Pucktrick} applied method type. The figure immediately reveals a decline in MCC between the two models across all experiments conducted. While CNN-LSTM shows a significant drop, TabNet exhibits only minimal declines}
    \label{fig:placeholder}
\end{figure}

\end{landscape}


At first glance this result might seem extremely unlikely, as we previously saw (in Section \ref{sec:res-B-single-label}) that corrupting \emph{just one} important feature allowed \emph{both} CNN-LSTM \emph{and} TabNet to have a significant drop in MCC; while in this other experiment (where that same feature is being corrupted alongside many other important features) the \textbf{TabNet MCC decay never accelerates into the super-additive collapse seen for the CNN-LSTM}. The TabNet progression is gently monotone in the noise level (e.g.\ the top-three group drifts $-0.004 \to -0.013 \to -0.039$ from $5\%$ to $50\%$ to $75\%$) and never accelerates into the super-additive collapse seen for the CNN-LSTM. 

Once again this big difference in between the two models can be explained by TabNet's per-decision-step sparse feature selection re-routes attention around the corrupted columns extremely effectively when several are hit at once as when only one is, so there is no compound cliff. This statement could further be backed up by taking a look at following graphs that show of the feature importance ranking changes, for each feature, alongside the multi-feature noise injection experiments:

\begin{itemize}
    \item For \textbf{TabNet}:
    \begin{itemize}
        \item For baseline's top features: we can clearly see that, as shown in image below, during the course of those experiments, the feature importance of those top 3 features drastically decreases. In fact, it can be seen that all three of what were once the most important features now rank fifth, sixth, and seventh out of the total of nine features (observing corruption level at 75\% of method \texttt{noise}): \begin{figure}[H]
            \centering
            \includegraphics[width=1\linewidth]{images//results/ExpB-multi-faeture-reRanking-top.png}
            \caption{Importance-rank drift of the corrupted feature in multifeature experiment. TabNet baseline top features re-rank across experiments. Graph show average rank above all experiments that used this column.}
            \label{fig:Importance-rank drift of the corrupted features in multifeature experiment. TabNet baseline top features re-rank across experiments.}
        \end{figure}
        
        \item For some of the baseline's lower tier features we can observe the exact opposite behavior: a significant increase in importance rank, signifying TabNet's ability to shift it's attention over other features: \begin{figure}[H]
            \centering
            \includegraphics[width=0.75\linewidth]{images//results/ExpB-multi-faeture-reRanking-bottom.png}
            \caption{Importance-rank drift of the corrupted feature in multifeature experiment. TabNet baseline bottom features re-rank across experiments. Graph show average rank above all experiments that used those columns.}
            \label{fig:Importance-rank drift of the corrupted feature in multifeature experiment. TabNet baseline bottom features re-rank across experiments.}
        \end{figure}
        Notably, one of the colums that significantly receives a rank boost even when being corrupted is \texttt{Fwd Seg Size Min}, which is the sole one that we identified as "Low-Fragility" inside the whole dataset inside Dataset Analysis Chapter \ref{sec:fragility}, once again solidifying this notion. 
    \end{itemize}

    \item For \textbf{CNN-LSTM}: 
    To support the claim that CNN-LSTM is not able to shift its attention to other features as effectively as TabNet (and consequently performs worse in these experiments as well), we can show the following graph illustrating the change in the ranking of all CNN-LSTM features across noise injection experiments, from which it is clear that no significant changes in rank occur, which inevitably leads the model to learn from corrupted columns:
    \begin{figure}[H]
        \centering
        \includegraphics[width=1\linewidth]{images//results/expB-cnn-lstm-multifeature-rerank.png}
        \caption{Importance-rank drift of the corrupted feature in multifeature experiment. All CNN-LSTM features re-rank across experiments. Graph show average rank above all experiments that used those columns.}
        \label{fig:Importance-rank drift of the corrupted feature in multifeature experiment. All CNN-LSTM features re-rank across experiments.}
    \end{figure}
    
\end{itemize}


\paragraph{Corruption on Multiclass-Task classification head:}
On the \textbf{multi-class} task the picture is subtler but equally telling.

For CNN-LSTM the \texttt{outliers} corruption of the six-feature top-and-middle group drives the CNN-LSTM multi-class MCC down by more than half its baseline at $75\%$ ($-0.526$), the single most severe feature-level degradation in the entire campaign; \texttt{noise} on the same group is comparably destructive ($-0.494$), whereas \texttt{missing} is markedly gentler ($-0.223$), because masked values are absorbed by the embedding/scaling pipeline rather than injected as out-of-distribution magnitudes. 

Meanwhile, TabNet pays a near-uniform penalty of $\approx -0.025$ to $-0.035$ that is already present at the lowest noise level and barely deepens as corruption intensifies (the top-three group moves only from $-0.028$ at $5\%$ to $-0.054$ at $75\%$). Decisively, this offset is essentially identical for the \emph{zero-variance control} group ($-0.028$ at $5\%$, $-0.039$ at $75\%$), whose features carry no information at all: the penalty therefore cannot be attributed to the loss of the corrupted signal, and is better read as a generic sensitivity of TabNet's minority-class resolution to \emph{any} perturbation of the training distribution (the same multi-class "tax" already noted in the single-feature regime, Section~\ref{sec:res-B-single-feat}), here confirmed to be independent of which, and how many, features are dirtied. The contrast with the CNN-LSTM is stark: on the worst group at $75\%$ TabNet's multi-class head loses $-0.051$ against the CNN-LSTM's $-0.526$, an order of magnitude less. 



% ------------------------------------------------------------
\subsubsection{Corruption of the target-label group}
\label{sec:res-B-multi-label}
% ------------------------------------------------------------

The seventh group corrupts the two targets \texttt{Label} and \texttt{label\_generic} \emph{jointly} with a single \texttt{labels} call. 

\begin{figure}[H]
    \centering
    \includegraphics[width=1\linewidth]{images//results/expB_multifeature_targets.png}
    \caption{Enter Caption}
    \label{fig:expB_multifeature_targets}
\end{figure}


The outcome is, at first sight, surprising: the catastrophic binary cliff of Section~\ref{sec:res-B-single-label} does \emph{not} reappear. Both models retain near-baseline binary MCC across the entire grid (TabNet $\Delta \approx -0.004$; CNN-LSTM between $+0.015$ and $-0.009$), while the multi-class task degrades only mildly (TabNet $-0.027$ to $-0.043$; CNN-LSTM $+0.033$ at $5\%$ down to $-0.086$ at $50\%$).

The explanation lies in the \emph{consistency} of joint corruption and is itself an informative result. When the two targets are corrupted together, a row whose multi-class label is reassigned from one attack type to another keeps its binary label malicious; reassignments among the many attack classes therefore overwhelmingly \emph{preserve} the benign/malicious distinction. The binary supervision is thus largely intact even at $50\%$ group corruption, which is why no inversion cliff forms, whereas the multi-class supervision is genuinely shuffled and slightly degrades. This stands in sharp contrast to the isolated binary-target flip, where every selected entry is guaranteed to cross the benign/malicious boundary. The lesson is methodologically important: \textbf{the impact of label noise depends not only on its quantity but on the precise mechanics of how it is realised}, whether corruption respects or violates the semantic coupling between related labels.

% ------------------------------------------------------------
\newpage
\subsection{Joint Considerations on Experiment B}
\label{sec: Joint Considerations on Experiment B}
% ------------------------------------------------------------

Taken together, the single-feature and multiple-feature scenarios of \red{Experiment~B} paint a coherent picture of how the two architectures respond to data-quality degradation.

\paragraph{Architectural inductive bias dictates robustness.}
TabNet's sparse sequential attention makes it the more robust model by a wide margin: it absorbs single- and even multi-feature corruption with negligible loss on the binary task, failing only when the supervision signal itself is inverted. CNN-LSTM is more sensitive to feature corruption but, by the same token, more \emph{plastic}: it is the only model capable of \emph{benefiting} from corruption, because it is the only one whose inductive bias (temporal windowing) admits a spurious shortcut that corruption can remove. Robustness and improvability are therefore two faces of the same architectural coin.

\paragraph{A controlled noise budget can be net-positive, but the window is narrow.}
The clearest cumulative break-even result is that, for the CNN-LSTM, there exists a corruption that is reliably and significantly \emph{beneficial} (\texttt{outliers}/\texttt{Timestamp}, monotone up to $+0.147$ binary and $+0.236$ multi-class), but it is highly specific: it requires corrupting a feature that encodes leakage, with a method that destroys that feature's structure, on a model that would otherwise overfit it. As soon as corruption is broadened to several informative features, or applied to a tabular model with no shortcut to remove, the net effect turns negative. The break-even point between regularisation and information destruction is thus real but tightly bounded, and locating it requires precisely the kind of dense, multi-seed sweep that the newly implemented Spark backend of PuckTrick (Chapter~\ref{chap:5-Empiric-analysis-models}) made computationally affordable.

\paragraph{The two tasks degrade at different rates.}
Across every scenario, the binary task is the more robust of the two and the multi-class task carries a persistent small penalty under almost any perturbation. The fine-grained resolution of rare attack subtypes is the first capability to erode and the last to recover, which is consistent with its dependence on minority-class signal that even modest corruption can wash out. Any deployment that cares about attack \emph{attribution} (not merely detection) must therefore treat data quality as a first-order concern.

\paragraph{Degradation is a readout of feature importance.}
Finally, the magnitude of degradation under corruption tracks the combined importance ranking with remarkable fidelity: \texttt{Dst Port} corruption hurts most, mid-tier features less, and zero-variance controls not at all. This means the noise-injection protocol doubles as an \emph{importance-attribution probe}: an independent, perturbation-based confirmation of the attention and permutation-based rankings of Section~\ref{sec:combined-ranking}.

% =======================================================================
\newpage
\section{Experiment A: Proof-of-Concept Results}
\label{sec:res-A}
% =======================================================================

\red{Experiment~A} is the smaller, proof-of-concept phase, conducted on the feature-pruned dataset over five seeds $S_A = \{1, 2, 42, 51, 84\}$, four corrupted columns (\texttt{FIN Flag Cnt}, \texttt{Down/Up Ratio}, \texttt{Protocol}, \texttt{Timestamp}) on dataset composed of 8 total features, and the seven-level grid $p_A$. As this experiment uses a different subset of features from full dataset, where as they were retained only high-signal features: all features are relevant for both labels (as explained before in Chapter \ref{sec:combined-ranking}). For this reason its clean baselines are stronger than \red{Experiment~B}'s on the reduced feature set:
\begin{itemize}
    \item for \textbf{TabNet}: $\mathrm{MCC}_{\text{bin}}/\mathrm{MCC}_{\text{mul}} = 0.896/0.886$
    \item fon \textbf{CNN-LSTM}: $\mathrm{MCC}_{\text{bin}}/\mathrm{MCC}_{\text{mul}} = 0.730/0.580$ 
\end{itemize}
Read in the light of \red{Experiment~B}, this phase is best understood as a concrete instance of attempting to re-create the regularization phenomenon characterized above, now visible across \emph{most} of the feature-set rather than concentrated on \texttt{Timestamp} alone. This is done in order to reach proof of concept for a real word application of both \texttt{Pucktick} library and this thesis theme of study: proving that deep learning NIDS can be improved with systematic noise injection. 

\paragraph{Disclaimer: he absence of formal significance}
No \red{Experiment~A} cell passes the $p < 0.05$ significance test; but this is a consequence of the test's statistical \emph{power} at five seeds, not of absent effects. As explained in Section~\ref{sec:res-protocol}, the significance screen is a paired Wilcoxon signed-rank test on the binary MCC, and with only $n = 5$ seeds this non-parametric test has a hard floor: even when all five seeds agree on the sign of the effect, the smallest two-sided $p$-value it can return is $\approx 0.0625 > 0.05$. No five-seed configuration can therefore ever be flagged, irrespective of how large or how tightly-bounded the effect is. And the effects \emph{are} often tightly bounded: the headline multi-class gains of Table~\ref{tab:res-A-cnnlstm} carry confidence half-widths as small as $\pm 0.014$ while being an order of magnitude larger, so they are \emph{descriptively} unambiguous even though the formal test must abstain (the exception is the \texttt{Protocol} control, whose $\approx\!\pm 0.5$ margins signal genuine seed disagreement and thus a truly inconclusive cell). The proper reading is that \red{Experiment~A} \emph{motivated} the hypothesis with large, mostly well-estimated effect sizes that its five-seed budget could not formally certify, and \red{Experiment~B}, with twenty seeds and the full dataset, then \emph{confirmed} the surviving component of that hypothesis (the CNN-LSTM/\texttt{Timestamp} regularization) at $p < 10^{-4}$ while delimiting its boundaries. The two phases are complementary: \red{Experiment B} proves whit certainty various effects of noise injections, while \red{Experiment A} can be considered as a real word application of \red{Experiment B} where we assume that the results have already been consolidated.

% -------------------------------------------------------------
\begin{landscape}
\subsection{Full Experimental grid}
% -------------------------------------------------------------

Instead of applying noise to all features presents in this Experiment sub-set (as previously done in \red{Experiment B}), noise was applied only on a smaller sub-set of 4 columns drawn from the reduced feature set, chosen to span distinct feature types and fragility levels:
$$\texttt{FIN Flag Cnt}\text{ (binary), }\texttt{Down/Up Ratio}\text{ (continuous), }\texttt{Protocol}\text{ (categorical), }\texttt{Timestamp}\text{ (date-time)} $$
Full experimental grid of all experiments is visible below in Images \ref{fig:ExpA_Cnn-lstm_full_grid_single_feature} and \ref{fig:CNN-ExpA_Tabnet_full_grid_single_feature}:
\begin{figure}[H]
    \centering
    \makebox[\textwidth][c]{\includegraphics[width=1\linewidth]{images/results/ExpA_fullGrid_CNN-LSTN.png}}
    \caption{Full result grid spectrum over MCC for CNN-LSTM in both prediction heads (binary task top, multitask task bottom)}
    \label{fig:ExpA_Cnn-lstm_full_grid_single_feature}
\end{figure}
\end{landscape}

\begin{landscape}
    \begin{figure}[H]
    \centering
    \makebox[\textwidth][c]{\includegraphics[width=1.3\linewidth]{images/results/ExpA_fullGrid_TabNet.png}}
    \caption{Full result grid spectrum over MCC for TabNet in both prediction heads (binary task top, multitask task bottom)}
    \label{fig:CNN-ExpA_Tabnet_full_grid_single_feature}
\end{figure}
\end{landscape}


\paragraph{For CNN-LSTM the proof of concept holds: controlled noise improves the weaker model.}
The headline result is that, for the CNN-LSTM, noise injection \emph{systematically} raises held-out MCC, with by far the largest and most consistent gains on the multi-class task: 

\begin{table}[H]
\centering
\caption{\red{Experiment~A}: CNN-LSTM multi-class $\Delta\mathrm{MCC} \pm \delta$ under \texttt{outliers} corruption (baseline $0.580$), where $\delta$ is the Student-$t$ $95\%$ half-width over the five seeds. Gains are large and \emph{tightly bounded} for three of the four features (margins as small as $\pm 0.008$); \texttt{Protocol} is the exception (see text), its $\approx\!\pm 0.5$ margins revealing that its five seeds disagree even on the sign. A margin of $\pm 0.000$ denotes a cell in which all five seeds returned a numerically identical MCC.}
\label{tab:res-A-cnnlstm}
\small
\resizebox{\linewidth}{!}{%
\begin{tabular}{lrrrrrrr}
\toprule
\textbf{Feature} & \textbf{1\%} & \textbf{5\%} & \textbf{10\%} & \textbf{20\%} & \textbf{30\%} & \textbf{50\%} & \textbf{75\%} \\
\midrule
\texttt{Timestamp}     & $+0.171{\pm}0.014$ & $+0.241{\pm}0.052$ & $+0.291{\pm}0.019$ & $+0.326{\pm}0.014$ & $+0.336{\pm}0.022$ & $+0.353{\pm}0.014$ & $+0.310{\pm}0.021$ \\
\texttt{FIN Flag Cnt}  & $+0.161{\pm}0.024$ & $+0.180{\pm}0.023$ & $+0.175{\pm}0.029$ & $+0.183{\pm}0.008$ & $+0.172{\pm}0.016$ & $+0.156{\pm}0.000$ & $+0.122{\pm}0.000$ \\
\texttt{Down/Up Ratio} & $+0.168{\pm}0.033$ & $+0.168{\pm}0.020$ & $+0.175{\pm}0.026$ & $+0.172{\pm}0.046$ & $+0.123{\pm}0.050$ & $+0.160{\pm}0.053$ & $+0.169{\pm}0.033$ \\
\texttt{Protocol}      & $-0.138{\pm}0.502$ & $-0.124{\pm}0.517$ & $-0.132{\pm}0.508$ & $-0.138{\pm}0.501$ & $-0.128{\pm}0.513$ & $+0.179{\pm}0.000$ & $+0.189{\pm}0.000$ \\
\bottomrule
\end{tabular}%
}
\end{table}

Table~\ref{tab:res-A-cnnlstm} reports the \textbf{CNN-LSTM multi-class improvement} under \texttt{outliers} corruption for the four features. The gains are substantial (up to $+0.353$ for \texttt{Timestamp}), positive across nearly the entire grid, and present for \texttt{missing} and \texttt{noise} as well, not only \texttt{outliers}. 

Yet, as we specified before for For the CNN-LSTM (in Section \ref{sec: Joint Considerations on Experiment B}), corruption of \texttt{Timestamp} using method \texttt{outliers} is highly specific for CNN-LSTM: it requires corrupting a feature that encodes leakage, with a method that destroys that feature’s structure, on a model that would otherwise overfit it. Therefore this doesn't represent a real-word scenario in it's fullest, as one would assume that some regularization techniques that help shuffle the data would have been used \emph{before} relying on noise injection to solve this very specific problem. Nonetheless this is still a solution to this hyper specific problem, that could be applied in very specific cases. 

Yet again, if we examine the results in Table \ref{tab:res-A-cnnlstm}, we can see that when using the \texttt{outliers} method, Timestamp is not the only feature that explicitly improves the MCC (in this case, for the second prediction head). This result can therefore guide us in finding a method \emph{consistent for both models}, which we will discuss in section \ref{parag: real-world application}

\paragraph{TabNet has no headroom and is unmoved.}
For TabNet, already near ceiling on the reduced set, the deltas are small and of mixed sign, with no systematic improvement. This is consistent with \red{Experiment~B}: a model that does not overfit a spurious signal has nothing to gain from corrupting it. The proof of concept is therefore specifically a statement about \emph{regularising an under-fit-prone temporal model}, not a universal claim about all architectures.


\paragraph{Real word applications.}
\label{parag: real word application}
AS we previously saw, \texttt{Pucktrick}'s \texttt{outliers} method was one that peaked our attention as it was the sole one that was capable of relayably increasing the MCC over the whole experimental grid (even if only on multiclass-head predictions). This suggests us that there is possibility of increasing models performance, but not by much. Or at least, that is what I would say if, in this particular case, we were not considering a scenario applicable to the real world. In this specific experiment, since this proof of concept is supposed to be based on real data (and not purely experimental data), it would not make sense to rely \emph{solely} on the trend (rise/fall) of the MCC for the two prediction models. In this particular case, we must also consider the \emph{meaning} of the binary predictions we make with our models. 
Given that we have:
\begin{itemize}
\tightlist
\item
  \textbf{Class 0 = benign packet}
\item
  \textbf{Class 1 = malicious packet}
\end{itemize}

We have that our predictions made with the models correspond to: 
\begin{itemize}
\item
  \textbf{True Positive (TP)} (\emph{true-positive prediction}): the
  package is actually malicious (1) and the model classifies it as
  malicious (1).
\item
  \textbf{True Negative (TN)} (\emph{true-negative prediction}): the
  package is actually benign (0) and the model classifies it as
  benign (0).
\end{itemize}

Whereas, \emph{on the other hand}, the error cases are:
\begin{itemize}
\item
  \textbf{False Positive (FP)}: The package is actually benign (0), but the model classifies it as malicious (1). This case corresponds to a \textbf{false alarm}. The NIDS reports an attack that does not actually exist.
\item
  \textbf{False Negative (FN)}: the packet is actually malicious (1), but the model classifies it as benign (0). This is a \textbf{undetected attack} (\emph{missed detection}). In other words, the NIDS allows malicious traffic to pass through, believing it to be legitimate. 
\end{itemize}

In a real-world scenario, a False-Negative is \textbf{a much more serious case} than a False-Positive. In the former, the entire corporate infrastructure is put at risk due to the failure to detect a malicious packet, whereas in the latter, only a tiny piece of information from a network flow is lost (which, incidentally, can be retrieved later). This information remained completely obscured in the cases of \red{Experiment B}, since these did not consider real-world scenarios but rather purely statistical analyses of the influence of \texttt{Pucktrick} on the dataset (where the dataset was considered purely as either tabular data or a time series, and nothing more).

For this reason, in this section we will focus on the insight provided by the MCC (namely, that using the \texttt{Outlier} method can yield a modest improvement in model performance over multiclass-head prediction), while also investigating the optimal form of noise injection based on the average rates of false positives and false negatives compared to the baseline models. We will then determine the best noise injection method among those examined by checking whether there is:
\begin{itemize}
    \item a decrease in false-negative cases: which implies a decrease in number of security breaches.
    \item an increase in false-positive cases: which implies an increase in number of false alarms rang.
\end{itemize}

Most interesting results from this research are reported in Image \ref{fig:ExpA_Cnn-lstm_confusion-matrix_single_feature} and Image \ref{fig:ExpA_tabnet_confusion-matrix_single_feature} below:

\begin{landscape}
\begin{figure}[H]
    \centering
    \makebox[\textwidth][c]{\includegraphics[width=1\linewidth]{images/results/ExpA_confusiuonMatrix_CNN-LSTM.png}}
    \caption{CNN-LSTM most enticing result that show clear reduction of average number of False-Negative (over 5 random-seeds) across the usage of \texttt{Outliers} applications}
    \label{fig:ExpA_Cnn-lstm_confusion-matrix_single_feature}
\end{figure}
\end{landscape}

\begin{landscape}
\begin{figure}[H]
    \centering
    \makebox[\textwidth][c]{\includegraphics[width=1\linewidth]{images/results/ExpA_confusiuonMatrix_TabNet.png}}
    \caption{TabNet most enticing result that show clear reduction of average number of False-Negative (over 5 random-seeds) across the usage of \texttt{Outliers} applications}
    \label{fig:ExpA_tabnet_confusion-matrix_single_feature}
\end{figure}
\end{landscape}

From the results shown in the figures, we can clearly see that for \emph{both models} it is possible to significantly reduce the number of false negatives. In fact, it is clearly visible:
\begin{itemize}
    \item \textbf{for CNN-LSTM}: all the changes shown in Figure \ref{fig:ExpA_Cnn-lstm_confusion-matrix_single_feature} lead \emph{also} to a very slight improvement in the MCC on the binary prediction head. The most striking case for this model remains the \texttt{Timestamp} feature, which manages to reduce the number of false negatives by as much as $1395$. However, even if we were to exclude this result due to the considerations we made previously (Section \ref{sec: Joint Considerations on Experiment B}), there would still be many other cases where the number of false negatives decreases drastically (effectively making the NIDS \emph{much more robust} compared to its other iterations). Striking examples include: 
    \begin{itemize}
        \item Use of \texttt{Outliers} on the ‘\texttt{FIN Flag Cnt}’ column: both the 20\% and 30\% thresholds yield excellent results, reducing the number of false positives by $261$ and $210$, respectively.
        \item Use of \texttt{Outliers} on the ‘\texttt{Protocol}’ column: both the 50\% and 75\% thresholds result in a decrease in the number of false negatives by $200$ and $485$ cases, respectively.
    \end{itemize}
    \item \textbf{For TabNet}: All the changes shown in Figure \ref{fig:ExpA_tabnet_confusion-matrix_single_feature} result in a slight (negligible) decrease in the MCC; however, as also observed for the CNN-LSTM model, the same features that led to a significant reduction in the number of false negatives also apply to this model. Specifically:
    \begin{itemize}
        \item use of \texttt{Outliers} on the ‘\texttt{FIN Flag Cnt}’ column: both the 50\% and 75\% applications lead to a reduction of approximately $\approx300$ false-negatives in both cases. 
        \item use of \texttt{Outliers} on the ‘\texttt{Protocol}’ column: for the single 50\% application, we see a sharp decrease in the number of false negatives by as many as $632$.
    \end{itemize}
\end{itemize}

An extremely interesting fact is that, for both models, the features corrupted using the \texttt{Outliers} method, that allow us to achieve a lower number of false negatives: \emph{coincide}.  However, we can take a note on the percentages used to achieve these improved results:
\begin{itemize}
    \item for the ‘\texttt{FIN Flag Cnt}’ column, the noise percentages \emph{do not match} between the two models.
    \item for the ‘\texttt{Protocol}’ column, the noise percentages \emph{match} between the two models, since both report that injecting 50\% noise into this column drastically reduces the number of FNs.
\end{itemize}
Taking this final information into account, we can therefore conclude that it is possible to systematically improve the performance of both models (and consequently of any deep learning-based NIDS built using \red{Experiment A}) by applying noise injection specifically to the ‘\texttt{Protocol}’ column.

We have thus concluded that by applying a thorough analysis of the dataset’s properties, it is possible to identify a feature (in our case, ‘\texttt{Protocol}’) which, when corrupted, \emph{consistently} (\emph{regardless of the model applied}) leads to a reduction in the occurrence of false positives: which, in this case, would correspond to a reduction in critical failures of the NIDS system. Although this may entail a small loss in overall accuracy (as occurred in our case with the loss in MCC value for the TabNet model), the net gain in real-world application is \emph{immense}, since it prevents a significant number of critical failures for the model. This leads us to conclude that, for a real-world application, it is possible to use noise injection tools such as \texttt{Pucktrick} to improve the performance of deep learning-based NIDS: thus, our proof-of-concept is validated. 


% =======================================================================
\newpage
\section{Conclusions}
\label{sec:conclusions}
% =======================================================================

This thesis set out to study both how noise injections affect learning on tabular and time-series data (major studies of \red{Experiment B}) and how controlled data-quality degradation affects deep-learning Network Intrusion Detection Systems (major studies of \red{Experiment A}). To this end the \texttt{PuckTrick} library, with the Spark backend developed in this work, was used to inject five families of noise (\texttt{missing}, \texttt{noise}, \texttt{outliers}, \texttt{duplicated} and \texttt{labels}) into the CIC-IDS2018 dataset, across a grid of nearly seven thousand experimental configurations and almost fourteen thousand individually trained models, on two complementary architectures: for tabular data TabNet, and for time-series data CNN-LSTM. The whole work rested first on an engineering foundation: the Spark re-engineering of the \texttt{PuckTrick} library, then unfolded into two experimental phases with distinct purposes: the analytical \red{Experiment~B} and the application-oriented \red{Experiment~A}. The conclusions are best stated against those three purposes, with the principal findings woven into the objective each one served.

% ------------------------------------------------------------
\subsection{Revisiting the Objectives of the Thesis}
\label{sec:obj-revisited}
% ------------------------------------------------------------

As said before, this thesis rested on three pillars, each pursuing a distinct objective: the re-engineering of the \texttt{PuckTrick} library with a Spark backend (Chapter~\ref{chap:oth}) and, sequentially, the two experimental campaigns it made possible (Chapter~\ref{chap:5-Empiric-analysis-models} and the present one). We can proudly say that all three objectives are met.


\paragraph{Objective library re-engineering: Rendering \texttt{PuckTrick} available on Spark: achieved.\newline}
The first pillar of this work, completed before any experiment could be run, was the extension of \texttt{PuckTrick} with an Apache Spark backend (Chapter~\ref{chap:oth}). The original Pandas implementation was built on row-level indexing and eager, in-memory operations. In order to make the newly implemented backend usable to the user, while not losing old Pandas functionalities, \texttt{Pucktrick} was re-architected into a three-tier, modular layout (\texttt{core/}, \texttt{pandas\_backend/}, \texttt{spark\_backend/} wich allowed to maintain old Pandas noise injection methods, while also adding the newly implemented methods in spark-backend) in which a \texttt{BackendFactory} transparently resolves the engine (automatically by dataset size, or on explicit demand) behind an \emph{unchanged} public API. All five error-injection modules (\texttt{duplicated}, \texttt{labels}, \texttt{missing}, \texttt{noisy}, \texttt{outliers}) were rewritten to respect Spark's columnar, lazily-evaluated paradigm, which required overcoming the absence of positional indices (through a surrogate \texttt{\_pucktrick\_row\_id}), the probabilistic nature of distributed sampling (through an oversample-and-truncate strategy that restores the \emph{exact} requested corruption percentage), and several DAG-materialization inefficiencies. Functional equivalence between the two backends was certified by a dedicated suite of unit and cross-backend integration tests, and the implementation was validated on a real five-node Standalone Spark cluster. The payoff is decisive and scalable. By comparing each Spark-\texttt{Pucktrick} method to its original Pandas-counterpart we can divide results into two branches:
\begin{itemize}
    \item Low Computational Complexity Methods: (\texttt{duplicated} and \texttt{missing}) those methods, due to implementing operations of inherently low computational complexity in Pandas-backend, they result very efficient even on large datasets. Yet, newly implemented spark counterparts still reaches reach break-even point on datasets larger than $\approx7$ million records, onwords.
    \item High Computational Complexity Methods: (\texttt{noise}, \texttt{outliers} and \texttt{labels}): those methods always result much faster than Pandas-counterpart, reaching break-even-point on all effectuated tests.
\end{itemize}
Another extremely important result is that by running the full five-method pipeline (running the five \texttt{Pucktrick} methods one after another), the Spark backend is already $8.2\times$ faster than Pandas at one million records and saturates at roughly $40$-$43\times$ beyond $32$ million. Crucially, that $\approx 8\times$ gain at the one-million-row scale is not incidental: it is precisely what rendered the $\approx 14$k-model campaign of the two experimental phases computationally tractable, binding the engineering contribution directly to the analytical one.

\paragraph{A reproducible analytical pipeline tied the two Experiments together. \newline}
Beyond the library development itself, a fully data-driven evaluation pipeline was built to recover every experimental coordinate (model, method, corrupted feature, noise level, seed) directly from artifact metadata rather than from hard-coded assumptions, so that a single re-run regenerates every table, figure and confidence interval of the results chapter. This is what allowed a grid of nearly seven thousand configurations to be navigated, cross-checked against its own negative controls, and reported with statistical rigor.


\paragraph{Objective of \red{Experiment~B}: a theoretical map of noise effects: achieved.\newline}
\red{Experiment~B} aimed not to improve a NIDS but to characterize, as completely as possible, the space of effects that controlled corruption produces on a tabular and a time-series classifier: These are the main case studies in this thesis; therefore, this is the most important experiment among all those conducted. In order to better understand full range of effects of noise injection, this experiment started from from the full-feature-dataset (79 features) and worked it's way down to a representative sub-set composed of nine features (plus a tenth one: \texttt{Timestamp}), which corresponds to the maximum computational limit that can be applied given the available computational resources, as the the full experimental grid expanded over 20 seeds, granting statistical relevance of achieved results (Student-$t$ intervals and paired Wilcoxon tests over 20 seeds), which meant requiring lots of time and computational resources. 
This experiment delivered on its promise of a comprehensive behavioural map, and the picture it produced is best read as a single connected story. Its first and most reassuring lesson is that both architectures are, on the whole, remarkably resilient to the corruption of a \emph{single} input feature: across the entire grid the binary MCC barely moves, the lone exception being the most discriminative column, \texttt{Dst Port} (most important column on CNN-LSTM and second most important column on TabNet), whose corruption degrades both models. This very regularity is itself the central diagnostic insight of the phase, because the magnitude of degradation tracks the combined feature-importance ranking so faithfully (maximal for \texttt{Dst Port}, negligible for the less important columns) that the noise-injection protocol effectively doubles as a model-agnostic \emph{importance probe}, independently corroborating the attention- and permutation-based rankings of Section~\ref{sec:combined-ranking}. Yet. against this backdrop of broad resilience, one corruption stands out as genuinely \emph{beneficial}: applying \texttt{outliers} to the \texttt{Timestamp} column raises the CNN-LSTM's MCC monotonically and significantly (up to $+0.147$ binary and $+0.236$ multi-class, $p<10^{-4}$), because permuting the training sequence before windowing strips away a spurious temporal-shortcut and thereby acts as a targeted \emph{regulariser} for model CNN-LSTM that would otherwise overfit on said temporal-shortcut. Tellingly, this gain is inexistent for the shortcut-free TabNet. This reveals noise injection to be not a generic booster but a \emph{precision instrument} for removing learnable leakage.

The same data also explain \emph{why} the two models diverge so sharply, showing that robustness is fundamentally an architectural property, legible even in the feature-importance dynamics: TabNet's sparse sequential attention re-routes around corrupted columns, demoting the dirtied features and \emph{promoting} clean ones (notably the low-fragility \texttt{Fwd Seg Size Min}). This outcome is further verified in experiments dedicated to corruption of multiple features. Corrupting even the six most important features at $75\%$ costs it only $\approx -0.05$ MCC: an order of magnitude less than the CNN-LSTM's super-additive collapse to $-0.526$ on the same group, the latter proved to being unable to re-rank and thus forced to keep learning from the corrupted signal. 

Most consequential of all, however, is the finding that corrupting the \emph{supervision signal} is categorically more dangerous than corrupting features: binary-label noise is harmless up to $\approx 50\%$ and then collapses off a cliff into outright inversion ($\Delta\mathrm{MCC}=-1.642$ for TabNet at $75\%$, a near-complete sign reversal), exactly as a two-state symmetric flip predicts. The multi-task probe moreover exposed two effects that a single-task study could never have seen: that label damage \emph{leaks} across heads, but only in the binary$\rightarrow$multi-class direction and only for the CNN-LSTM, whose shared context vector lets the binary collapse drag its multi-class head down by $-0.536$ while TabNet's attention-disentangled heads barely move; and that the harm depends on the \emph{mechanics} of corruption rather than its sheer quantity, since jointly corrupting both targets produced no binary cliff at all, the reassignment of an attack to another attack class preserving the benign/malicious distinction. 

Finally, across every scenario binary detection proved more robust than fine-grained attribution, the latter carrying a near-uniform $\approx -0.03$ penalty present already at the lowest noise level and even when only zero-variance features were touched, which marks it as a generic sensitivity of minority-class resolution rather than a loss of specific information. The resulting picture is coherent, statistically grounded and internally validated by its own negative controls; the analytical objective is therefore fully met.


\paragraph{Objective of \red{Experiment~A}: a real-world proof of concept: achieved. \newline}
\red{Experiment~A} asked the more pragmatic question of whether targeted noise injection could \emph{improve} a deployable deep NIDS. This experiment had a similar setting to \red{Experiment B}, yet over operating on the full dataset, this experiment started from an elaborated 44-features subset where all redundant features and zero-variance features were dropped. Though a similar analysis done in previous experiment then, only a fewer subset of 8-features statistically relevant for both models were retained (once again those restraints were due the computational prowess and long experimental time).
This experiment, too, delivered on its promise, but only once the right yardstick was adopted: and that shift of yardstick is itself one of the thesis's central lessons. The proper measure of a data-quality intervention for security is not the aggregate MCC but the \emph{asymmetric operational cost} of the errors a model makes, where a false negative (a missed attack) is far graver than a false positive (a false alarm). Read through that lens it was found possible to improve a deep NIDS in a crucially \emph{model-independent} way: corrupting the \texttt{Protocol} column with the \texttt{outliers} method at $50\%$ consistently reduced the number of false negatives for \emph{both} TabNet and CNN-LSTM (by $\approx 600$ and $\approx 200$ missed detections respectively) at the price of at most a negligible MCC loss. This manipulation, that the purely statistical reading of \red{Experiment~B} would have dismissed as unremarkable, yields a substantial real-world safety gain. A property of the dataset, identified through systematic analysis, can thus be deliberately exploited to make a deep NIDS materially safer in deployment, which validates the proof of concept: \texttt{PuckTrick}-style noise injection is a usable, practical tool for improving deep-learning intrusion detection, provided the target column, method and intensity are chosen with knowledge of both the data and the cost structure of the task.




\paragraph{}
In closing, this thesis has shown that controlled data corruption, far from being merely a threat to be defended against, can be turned into a deliberate, well-characterized \emph{tool}: one that probes what a model has learned, exposes the boundaries of its robustness, and, when wielded with knowledge of the data and the cost structure of the task, can make a deep-learning intrusion detection system measurably safer, thus confirming all of our premises.

-----------------------------------------------------------------------------------------


% ------------------------------------------------------------
\newpage
\subsection{Future Work}
\label{sec:future-work}
% ------------------------------------------------------------

Several directions follow naturally, on both the experimental and the library-engineering fronts.

\paragraph{On the experimental and modeling side.}
\begin{itemize}
    \item \textbf{Generalization beyond the present scope.} Finally, replicating the campaign on further architectures and datasets would establish how far the architecture-dependent conclusions of this thesis (attention-based robustness versus temporal-model plasticity) generalist across the broader landscape of deep-learning intrusion detection.
    
    \item \textbf{Extending the operational analysis to the full grid.} The false-negative / false-positive lens applied to \red{Experiment~A} should be extended to the complete feature set and to the multi-feature groups of \red{Experiment~B}, and ultimately validated in an online or streaming detection setting closer to real deployment.

    \item \textbf{Treating noise injection as a tunable hyperparameter.} Rather than discovering beneficial corruptions by exhaustive sweep, future work could \emph{search} directly for the optimal (feature, method, percentage) triple: optimizing explicitly for the operational objective (e.g.\ minimal false negatives at bounded false-positive cost) rather than for MCC, turning the proof of concept of \red{Experiment~A} into an automated data-centric tuning procedure.
    
    \item \textbf{Formalizing and mitigating multi-task contamination.} The asymmetric binary$\rightarrow$multi-class leakage observed in the CNN-LSTM invites a formal characterization of how a corrupted supervisory gradient propagates through a shared encoder, and motivates architectural countermeasures (such as attention-disentangled or gated task-specific heads) that could ''firewall'' one task from the catastrophic failure of another.    
\end{itemize}

\paragraph{On the library-engineering side.}
\begin{itemize}
    \item \textbf{Additional backends.} The \texttt{BaseBackend} abstraction was deliberately designed so that new engines can be added without touching the public API; natural candidates are a Polars or DuckDB backend for fast single-machine execution and a GPU/RAPIDS (\texttt{cuDF}) backend for very large in-memory workloads, each slotting in beneath the existing \texttt{BackendFactory}.
    
    \item \textbf{Closing the residual performance gaps.} Two specific weaknesses identified in the benchmarks deserve attention: lowering the $\approx 7$-million-row break-even point of the low-complexity \texttt{duplicated} and \texttt{missing} methods by further trimming the fixed DAG-construction and Python--JVM serialisation overhead, and fully porting the date-column \texttt{noise} path to native Spark date handling.

\end{itemize}

\end{document}

